% Main survey paper — AASTeX 6.3.1 (AJ/ApJ; journal chosen at submission).
% Figures: copy the referenced PNGs from results/ next to this file before
% compiling/submitting. Machine-readable tables: paper/table1_dispositions.csv
% and paper/table2_planet_limits.csv.
% Open items are marked [PENDING]/[TODO] in the text.
\documentclass[twocolumn]{aastex631}

\newcommand{\msini}{m\sin i}
\newcommand{\mjup}{M_{\rm Jup}}
\newcommand{\mearth}{M_\oplus}

\begin{document}

\title{Dynamics as a Scale: A Homogeneous Survey of Stability-Derived
True-Mass Limits for Non-Transiting Radial-Velocity Exoplanets}

\author{Shaurya Ranjan}
\affiliation{Independent Researcher}
\email{shauryaranjan911@gmail.com}

\begin{abstract}
Radial-velocity detections constrain only the minimum mass $\msini$ of a
planet; the true mass can exceed it by an arbitrary factor for face-on
orbits. In multi-planet systems, however, long-term dynamical stability
bounds that factor: heavier, more face-on configurations perturb their
neighbors strongly enough to self-destruct within the system's lifetime.
Discovery papers occasionally apply this constraint to individual systems
with heterogeneous methods; no uniform survey exists. We present a
homogeneous $N$-body stability survey of 56 multi-planet RV systems drawn
from the NASA Exoplanet Archive, sweeping each system's line-of-sight
inclination and sampling all orbital parameters within their published
uncertainties. We validate the pipeline against GJ~876, recovering the
literature stability constraint ($i > 25\arcdeg$ vs.\ published
$i \gtrsim 20\arcdeg$). We report 95\% credible true-mass upper limits
under an isotropic orientation prior for 123 planets in 53 systems: 61
planets in 23 systems are tightened below 3.0 times their minimum masses
(median limit $2.68\times$ among them), and 25 giant planets are
certified below the deuterium-burning limit by dynamics alone. The pipeline additionally documents, and guards
against, three classes of catalog-parameter pitfalls that dynamical
consistency exposes, including published eccentricity point-estimates
incompatible with the survival of their own systems.
\end{abstract}

% TODO: verify UAT concept numbers at submission
\keywords{Exoplanet dynamics --- Exoplanet catalogs --- N-body
simulations --- Radial velocity}

\section{Introduction} \label{sec:intro}

The radial-velocity method measures a planet's gravitational reflex along
the line of sight only, leaving the orbital inclination $i$ unconstrained
and delivering the minimum mass $\msini$ rather than the mass. For an
isotropic distribution of orientations the correction is usually modest
--- the median of $1/\sin i$ is 1.15, and 95\% of orientations imply a
true mass below $3.2\,\msini$ --- but for any individual planet the
catalog value is a floor, not a measurement. The consequences propagate
into everything built on RV masses: mass functions and occurrence rates
inherit a one-sided bias, the planet/brown-dwarf boundary is porous for
every giant whose $1/\sin i$ could exceed $\sim$$13\,\mjup/(\msini)$, and
models of individual systems' formation and composition rest on lower
limits treated, in practice, as values.

In multi-planet systems, dynamics supplies the missing upper bound. A
coplanar system tilted toward face-on hides planets that are heavier by
$1/\sin i$; heavier planets perturb one another more strongly; and
sufficiently heavy versions destabilize on timescales far shorter than
the systems' ages. Configurations that cannot survive to the present are
excluded, so the survival requirement converts each system's continued
existence into a ceiling on its planets' true masses. The idea is not
new: discovery papers have long applied stability cuts to individual
systems (GJ~876 being the canonical example; e.g.,
\citealt{Gozdziewski2002} for 47~UMa; \citealt{Vogt2014} for HD~141399;
\citealt{Basant2025} for Barnard's star), and
\citet{Tamayo2021} developed stability-constrained characterization into
a general method, applied to the transiting system Kepler-23. What does
not exist is uniformity. \citet{LaskarPetit2017} classified 131 systems
by AMD-stability but did not invert stability into per-planet mass
limits; \citet{Volpi2019} constrained mutual inclinations of ten RV
systems with first-order secular theory; \citet{Qin2025} propose
breaking the degeneracy through von Zeipel--Lidov--Kozai cycle
monitoring in compact binaries; the remaining literature is a
scatter of single-system analyses with mutually incompatible integrators,
instability criteria, timescales, and statistical definitions. No two of
their mass limits are comparable, and no population statement can be
assembled from them.

This paper presents the first homogeneous survey: one pipeline, one
instability criterion, one statistical definition, applied to every
qualifying multi-planet RV system in the NASA Exoplanet Archive, with
orbital parameters sampled within their published uncertainties and the
machinery validated against systems whose stability constraints are
independently known. We report 95\% credible true-mass limits under an
isotropic orientation prior for the planets of 56 systems, a byproduct
catalog of parameter-consistency flags, and --- as a methodological
contribution in its own right --- three classes of catalog-parameter
pitfalls that dynamical auditing exposes (\S\ref{sec:artifacts}).

\section{Sample Selection} \label{sec:sample}

We select systems from the NASA Exoplanet Archive's composite parameters
table (\texttt{pscomppars}, accessed 2026 July 9--13) requiring at least
two planets, periods and masses for all of them, a stellar mass, and at
least one non-transiting planet whose mass provenance is $\msini$ --- the
planets the survey can constrain. Because dynamical interactions weaken
steeply with separation, we further require at least one adjacent pair
with a period ratio below 6. This yields 148 qualifying systems.

Each system is ranked by an analytic pre-estimate of where its ceiling
should lie: the swept mass factor $f_{\rm crit}$ at which the most
fragile adjacent pair, with masses scaled by $f$ and separations measured
in mutual Hill radii (eccentricity-reduced), reaches the two-planet
overlap criterion $\Delta = 2\sqrt{3}$ \citep{Gladman1993}. Multi-planet
systems destabilize at wider spacings, so $f_{\rm crit}$ is conservative.
The survey band takes $f_{\rm crit} \in (1.05, 60)$ with
$e_{\rm max} \le 0.75$ and $P_{\rm inner} \ge 1.5$~d --- 56 systems. The
excluded $f_{\rm crit} \approx 1$ block is dominated by
resonance-protected pairs whose survival depends on phase coherence that
composite catalog parameters cannot supply
(\S\ref{sec:artifacts}, \S\ref{sec:discussion}); it is deferred rather
than constrained. Table~\ref{tab:dispositions} gives the full disposition
of all 148 systems: 53 surveyed, 3 flagged, and 92 excluded (42
resonance-protected, 48 too weakly interacting for dynamics to bite at
any allowed mass, and 9 with sub-1.5-day inner periods, categories
overlapping).

\section{Methods} \label{sec:methods}

\subsection{Configuration Sampling} \label{sec:sampling}

A coplanar system tilted rigidly with respect to the line of sight leaves
its internal dynamics unchanged while every $\msini$ mass scales as
$1/\sin i$; we therefore sweep $i$ over a 17-point grid from $90\arcdeg$
(edge-on, masses at their floors) to $5\arcdeg$ (masses at $11.5\times$
the floors), holding planets with transiting or astrometric true masses
fixed. At each grid point we integrate 20 independent realizations,
drawing the stellar mass and every planet's mass, period, eccentricity,
and periastron longitude from truncated normal distributions built on the
symmetrized catalog $1\sigma$ uncertainties (point values where errors
are unreported; eccentricities truncated to $[0, 0.95]$). Unknown angles
are drawn uniformly. Orbital phases are randomized for composite-catalog
systems and taken from the fit's time of periastron only for
single-reference parameter sets, for the reasons documented in
\S\ref{sec:artifacts}.

\subsection{Integration and Instability Criterion} \label{sec:integration}

Realizations are integrated with \textsc{rebound}'s WHFast
\citep{ReinLiu2012, ReinTamayo2015} with timestep
$dt = \min_p P_p (1-e_p)^{3/2} / 20$, resolving the fastest perihelion
passage. Instability is a close encounter within the sum of Hill radii
(particles carry their Hill radius as collision radius) or escape beyond
20 times the widest semi-major axis. Integration time is uniform in
inner-planet orbits --- $10^5$ and $10^6$ for the two survey depths ---
making the effort per system architecture-independent; convergence
between the depths is examined in \S\ref{sec:depth}. Sensitivity of the
final limits to these choices is quantified on a six-system
representative subset spanning the survey's architectures: halving $dt$
changes 95\% credible limits by $\le 0.03$ in the mass factor; collision
radii scaled to $0.5\times$ and $0.25\times$ the Hill radius change them
by $\le 0.08$; replacing exact coplanarity with mutual inclinations drawn
from Rayleigh distributions of $\sigma = 2\arcdeg$ and $5\arcdeg$ changes
them by $\le 0.4$. None approaches the survey's grid resolution in $i$.

\subsection{Statistical Definitions} \label{sec:stats}

Each system yields a survival curve $S(i)$, the fraction of realizations
surviving to $t_{\rm max}$ at inclination $i$. Because parameter draws
can be fragile at any tilt, constraints are defined relative to the
edge-on baseline: systems with $S(90\arcdeg) < 0.5$ are reported as
parameter-consistency flags --- statements about the published
parameters, not mass limits. For unflagged systems the primary statistic
is the 95\% credible limit on the mass factor under an isotropic
orientation prior (uniform in $\cos i$), with posterior density
proportional to $S(i)$: the factor below which 95\% of the surviving
posterior lies. With no dynamical information this limit is 3.20 for
every RV planet; the survey's information content is the tightening below
that value. A secondary threshold ceiling (smallest factor with
$S < 0.5 \times S(90\arcdeg)$) is tabulated for comparison with prior
single-system work.

\subsection{Catalog-Parameter Pitfalls Guarded Against}
\label{sec:artifacts}

Three recurring pitfalls in catalog parameters, each encountered because
it broke an analysis during survey development, are guarded against as
follows.

\begin{enumerate}
\item Eccentricity point-estimates dynamically inconsistent with system
survival (Barnard's star: $e = 0.03$--$0.08$ kills the system in kyr;
$e$ drawn within errors resolves it --- independently established by
\citealt{Basant2025}, whose fits favor $e < 0.02$). Guard: sample within
uncertainties, never integrate point values alone (\S\ref{sec:sampling}).
\item Composite tables (\texttt{pscomppars}) mix parameters across
incompatible fits --- fatal for phase-coherent (resonant)
initialization. Guard: randomize phases for composite parameters; use
fit phases only from single-reference sets (\S\ref{sec:sampling}).
\item Fit epochs stored in the time-of-periastron column
\citep[GJ~876;][]{Rivera2010}: read naively, the column implies all
planets at periastron simultaneously. Guard: reject identical periastron
times across a system before using any phase information.
\end{enumerate}

\section{Validation} \label{sec:validation}

The pipeline is validated against two systems whose stability constraints
are independently known, chosen to bracket the two dynamical regimes the
survey encounters: a resonance-protected system whose survival depends on
its orbital phases (GJ~876) and a near-resonant system that is stable
regardless of them (HD~45364). A third, external cross-check appears in
\S\ref{sec:catalog}: for Barnard's star, the survey's ceilings agree with
the independent stability-derived masses of \citet{Byrne2026}.

\subsection{GJ 876: The Phase-Protected Anchor} \label{sec:gj876}

GJ~876 is the canonical stability-constrained system. The Archive's
default solution is the coplanar dynamical fit of \citet{Rivera2010} at
$i = 59\arcdeg$: true masses, with the phases of the 4:2:1 Laplace
resonance at their libration center. We reconstruct the equivalent
minimum-mass system by scaling the fitted masses by $\sin 59\arcdeg$, fix
the mean anomalies at the fit epoch to the published values (their
Table~3 --- the Archive's time-of-periastron column stores the epoch
itself for this reference, artifact Class~3 of \S\ref{sec:artifacts}, so
phases cannot be derived from it), and sweep a 13-point inclination grid
with phases held fixed (5 draws per point, 5{,}000~yr per realization
$\approx 10^6$ inner-planet orbits).

The validation makes three independent demands
(Figure~\ref{fig:validation}). First, the published configuration,
entering the sweep at mass factor $1/\sin 59\arcdeg = 1.17$, must
survive: it does. Second, the recovered constraint must match the
literature: survival is zero at every grid inclination at or below
$25\arcdeg$ and recovers above it, giving $i > 25\arcdeg$
($m < 2.37\times\msini$) against the published $i \gtrsim 20\arcdeg$ ---
consistent, and conservative in the safe direction. Third, the same
masses with randomized phases must not reliably survive, or phase
protection would be irrelevant to the problem: survival collapses to
12\%. The resonance protection is real, and the phase-constrained
initialization captures it.

\subsection{HD 45364: The Phase-Agnostic Anchor} \label{sec:hd45364}

HD~45364 anchors the opposite regime. Its two giants sit in a 3:2
commensurability with low eccentricities \citep{Correia2009}, and the
\citet{Li2022} solution is stable without phase assistance: initialized
with the fit's own phases, the published configuration survives the full
$10^5$~yr integration, and randomizing the phases leaves survival at
100\% --- the constraint owes nothing to resonance protection. The
survival cliff is correspondingly smooth rather than resonance-sharpened:
complete survival for $i \ge 45\arcdeg$, declining to zero by
$i \approx 8.5\arcdeg$ (17-point grid, 10 draws per point), integrating
under the isotropic prior to a 95\% credible limit of
$m < 2.29\times\msini$. Together the anchors bracket the survey
population --- one ceiling that exists because of its phases, one
indifferent to them --- and the same pipeline recovers both with no
per-system tuning.

\subsection{Depth Convergence} \label{sec:depth}

A survey whose limits depended on integration length would be reporting
its clock, not the dynamics. Every system therefore runs at two depths,
$10^5$ and $10^6$ inner orbits, with identical parameter draws (seeds
keyed to the system name), and the threshold ceilings are compared.
Deeper integration can only reveal instability, never undo it, so
ceilings must move one way --- and they do: among the 53 unflagged
systems (flagged systems excluded at either depth), 9 ceilings tightened,
5 held exactly, none loosened, and 7 systems acquired a ceiling only at
$10^6$ (Figure~\ref{fig:depth}). This is the third consecutive survey
generation with strictly monotone behavior, and it is what makes the
reported limits conservative: a longer clock can only strengthen them
(\S\ref{sec:discussion}). \textbf{[PENDING: $10^7$-orbit extension of the
16 non-converged systems --- run launched 2026-07-14.]}

\section{Results} \label{sec:results}

\subsection{The Catalog} \label{sec:catalog}

Of the 56 surveyed systems, 53 pass the edge-on consistency test and
yield 95\% credible mass limits for 123 non-transiting planets
(Table~\ref{tab:limits}; Figure~\ref{fig:dumbbells}). Sixty-one planets
in 23 systems are dynamically tightened below $3.0\times$ their minimum
masses --- meaningfully inside the $3.20\times$ that isotropy alone
guarantees --- with a median limit of $2.68\times$ among them.
Twenty-five giants whose face-on orientations would have exceeded the
deuterium-burning limit are certified planetary at 95\% credibility.
Individual highlights: Barnard's star's four sub-Earths are capped at
$2.18\times$ their minimum masses (0.42--0.73\,$\mearth$) --- consistent
with, and independently confirming, the stability-derived masses of
0.19--0.84\,$\mearth$ in \citet{Byrne2026}; 47~UMa at $2.35\times$;
HD~141399's four giants at $2.71\times$; and K2-18~c, the non-transiting
companion of the habitable-zone sub-Neptune K2-18~b, at
$2.87\times \approx 21.6\,\mearth$, excluding a giant-planet identity.
(The isotropic prior is conservative for K2-18~c: because K2-18~b
transits, only a small mutual inclination is required for c to avoid
transiting, so near-edge-on orientations are likely a priori;
\citealt{Cloutier2017}.) The tightest factor in
the catalog is a caution against conflating the two headline statistics:
TYC~1422-614-1's pair is capped at $2.01\times$ --- the survey's
strongest relative constraint --- yet planet~c's minimum mass of
$10\,\mjup$ puts even that ceiling at $\approx$$20\,\mjup$, above the
deuterium-burning limit. A tight mass factor and a certified planetary
identity are independent properties.

\subsection{Population Statement} \label{sec:population}

The cumulative distribution of credible limits (Figure~\ref{fig:cdf})
answers the survey's title question: exactly half the qualifying planets
(61 of 123) have their permitted mass range tightened beyond geometry,
the tightest fifth to below $\sim$$2.6\times$, and no limit exceeds the
isotropic bound beyond the inclination grid's discretization (the single
nominal exception, TOI-1062~c at $3.22\times$ against the analytic 3.20,
is an unconstrained posterior rendered on the 17-point grid) --- dynamics
only ever removes orientations. Read as demographics: for packed
multi-planet RV systems, the ``hidden monster'' tail of the minimum-mass
distribution is largely a geometric fiction; the surviving configurations
concentrate near edge-on masses.

\subsection{Parameter-Fragility Flags} \label{sec:flags}

Three systems fail the edge-on consistency test in more than half their
parameter draws: GJ~667~C, HIP~57274, and Kepler-20. These are
data-quality results, not mass limits: the published parameter sets, with
their stated uncertainties, do not support long-lived configurations at
any tilt. For GJ~667~C this echoes the long-standing dispute over its
planet multiplicity; for the others it flags eccentricity solutions that
their own dynamics reject. A joint dynamical--RV refit is the indicated
follow-up in all three cases.

\section{Discussion} \label{sec:discussion}

Three limitations bound the interpretation of the catalog. First, the
sweep is coplanar: the system is tilted as a rigid unit, so mutual
inclinations enter only through the sensitivity variants
(\S\ref{sec:integration}), which shift the credible limits by at most
$\sim$0.4 in the mass factor for dispersions up to $5\arcdeg$. Strongly
non-coplanar architectures --- precisely those the von
Zeipel--Lidov--Kozai literature targets --- are outside this survey's
assumptions and their limits should be read as conditional on approximate
coplanarity. Second, $10^6$ inner orbits is short against Gyr ages. Every
depth comparison we have run --- three survey generations --- moved
limits monotonically tighter or not at all, so the reported ceilings are
conservative: longer clocks can only strengthen them. The converse claim
(that a surviving configuration is stable forever) is never made. Third,
the inputs are catalog parameters, sampled within symmetrized published
uncertainties rather than refit from radial velocities. The
parameter-fragility flags quantify where this matters most; for the
flagged systems the correct next step is a joint refit, not a mass limit.

The excluded resonance-protected block (adjacent period ratios within a
few percent of small-integer commensurabilities, screened out at
selection) is addressable with the machinery validated on GJ~876:
single-reference parameter sets carrying fit phases, initialized at the
observed libration. Extending the survey across that block would add some
of the most tightly packed --- and therefore most constrainable ---
systems in the catalog.

The catalog has immediate uses beyond its headline statement. Planets
whose dynamical ceiling sits close to their $\msini$ floor are the best
targets for Gaia astrometric confirmation, since their permitted mass
range is narrowest; planets whose ceiling crosses the deuterium-burning
limit are certified planetary without further observation, cleaning
brown-dwarf contamination out of the giant-planet mass function; and the
parameter-fragility flags hand RV teams a prioritized list of systems
whose published solutions their own data no longer support.

\section{Data Availability} \label{sec:data}

Code and per-run records are available at
\url{https://github.com/Shaun2highh/dynamass}. Catalogs and per-run
records will be archived at Zenodo \textbf{[TODO: DOI --- mint after
RNAAS note acceptance]}. This work used the NASA Exoplanet Archive
\texttt{pscomppars}
\dataset[(DOI: 10.26133/NEA13)]{https://doi.org/10.26133/NEA13} and
\texttt{ps} \dataset[(DOI: 10.26133/NEA12)]{https://doi.org/10.26133/NEA12}
tables, accessed 2026 July 9--13; we thank the Archive's maintainers,
whose stewardship makes surveys like this possible.

Artificial-intelligence assistance (Claude, Anthropic) was used for code
development and manuscript drafting; all results were independently
verified and are the author's responsibility.

\software{\textsc{rebound} \citep{ReinLiu2012}, WHFast
\citep{ReinTamayo2015}, \textsc{numpy}, \textsc{pandas},
\textsc{matplotlib}}

% ---------------------------------------------------------------- figures

\begin{figure*}[ht!]
\centering
\includegraphics[width=0.48\textwidth]{gj_876_validation.png}
\includegraphics[width=0.48\textwidth]{hd_45364_validation.png}
\caption{Validation anchors. Left: GJ~876 survival vs.\ inclination with
the \citet{Rivera2010} fit phases held fixed; survival is zero at
$i \le 25\arcdeg$, recovering the literature constraint. Right:
HD~45364 \citep{Li2022}: a smooth, phase-agnostic survival cliff
(complete survival at $i \ge 45\arcdeg$, zero by $i \approx 8.5\arcdeg$).}
\label{fig:validation}
\end{figure*}

\begin{figure*}[ht!]
\centering
\includegraphics[width=0.9\textwidth]{floor_vs_ceiling_final.png}
\caption{Mass floors ($\msini$) and 95\% credible dynamical ceilings for
the 123 constrained planets of the 53 unflagged systems, sorted by
credible mass factor. Flagged systems are excluded.}
\label{fig:dumbbells}
\end{figure*}

\begin{figure}[ht!]
\centering
\includegraphics[width=\columnwidth]{depth_convergence_prod.png}
\caption{Depth convergence: per-system threshold ceilings at $10^5$ vs.\
$10^6$ inner-planet orbits with identical parameter draws. Ceilings only
tighten (9) or hold (5); 7 systems acquire a ceiling only at depth; none
loosen.}
\label{fig:depth}
\end{figure}

\begin{figure}[ht!]
\centering
\includegraphics[width=\columnwidth]{m95_cdf_prod.png}
\caption{Cumulative distribution of the 123 95\% credible mass-factor
limits. The dashed line marks the isotropic-prior bound (3.20) that
geometry alone guarantees; half the sample is tightened inside it.}
\label{fig:cdf}
\end{figure}

% ----------------------------------------------------------------- tables

\begin{deluxetable}{llrrr}
\tablecaption{System dispositions (excerpt)\label{tab:dispositions}}
\tablehead{
\colhead{System} & \colhead{Disposition} & \colhead{$f_{\rm crit}$} &
\colhead{$S(90\arcdeg)$} & \colhead{$m_{95}$ factor}
}
\startdata
Barnard's star & surveyed & 1.60  & 0.55 & 2.18 \\
47 UMa         & surveyed & 1.34  & 0.55 & 2.35 \\
HD 141399      & surveyed & 3.05  & 0.85 & 2.71 \\
K2-18          & surveyed & 17.92 & 1.00 & 2.87 \\
GJ 667 C       & flagged  & 5.71  & 0.15 & \nodata \\
HIP 57274      & flagged  & 1.23  & 0.45 & \nodata \\
Kepler-20      & flagged  & 4.25  & 0.40 & \nodata \\
24 Sex         & excluded\tablenotemark{a} & 1.00 & \nodata & \nodata \\
TOI-500        & excluded\tablenotemark{b} & 32.86 & \nodata & \nodata \\
\enddata
\tablenotetext{a}{Resonance-protected / at the stability limit
($f_{\rm crit} \approx 1$).}
\tablenotetext{b}{Inner period below 1.5~d.}
\tablecomments{Disposition of all 148 qualifying systems: 53 surveyed, 3
flagged (edge-on survival $< 0.5$), 92 excluded with reasons
(resonance-protected, weakly interacting, or sub-1.5-day inner period).
The full machine-readable table is available as
\texttt{table1\_dispositions.csv}, regenerated by
\texttt{scripts/make\_table1.py}.}
\end{deluxetable}

\begin{deluxetable}{llrrr}
\tablecaption{Per-planet 95\% credible true-mass limits (ten tightest)
\label{tab:limits}}
\tablehead{
\colhead{System} & \colhead{Planet} & \colhead{$\msini$ ($\mearth$)} &
\colhead{$m_{95}$ factor} & \colhead{$m_{95}$ ($\mearth$)}
}
\startdata
TYC 1422-614-1 & b & 794.5 & 2.01 & 1600 \\
TYC 1422-614-1 & c & 3178  & 2.01 & 6400 \\
Barnard's star & b & 0.299 & 2.18 & 0.65 \\
Barnard's star & c & 0.335 & 2.18 & 0.73 \\
Barnard's star & d & 0.263 & 2.18 & 0.57 \\
Barnard's star & e & 0.193 & 2.18 & 0.42 \\
HD 95089       & b & 400.4 & 2.31 & 923  \\
HD 95089       & c & 1096  & 2.31 & 2527 \\
HIP 111909     & b & 384.5 & 2.33 & 897  \\
HIP 111909     & c & 257.4 & 2.33 & 601  \\
\enddata
\tablecomments{The ten tightest of the 123 constrained planets, sorted by
95\% credible mass factor. The full machine-readable table
(\texttt{table2\_planet\_limits.csv}, regenerated by
\texttt{scripts/make\_table2.py}) additionally lists masses in $\mjup$,
the secondary threshold ceiling, and the inclination floor.}
\end{deluxetable}

% ------------------------------------------------------------- references

\begin{thebibliography}{}
% TODO: fill ApJL volume/page for Basant2025 and Qin2025 at submission
\bibitem[Basant et al.(2025)]{Basant2025} Basant, R., Luque, R., Bean,
J.~L., et al. 2025, ApJL, doi:10.3847/2041-8213/adb8d5
\bibitem[Byrne et al.(2026)]{Byrne2026} Byrne, X., Guimond, C.~M.,
Bonsor, A., et al. 2026, MNRAS, accepted (arXiv:2606.23310)
\bibitem[Cloutier et al.(2017)]{Cloutier2017} Cloutier, R.,
Astudillo-Defru, N., Doyon, R., et al. 2017, A\&A, 608, A35
\bibitem[Correia et al.(2009)]{Correia2009} Correia, A.~C.~M., Udry, S.,
Mayor, M., et al. 2009, A\&A, 496, 521
\bibitem[Go\'zdziewski(2002)]{Gozdziewski2002} Go\'zdziewski, K. 2002,
A\&A, 393, 997
\bibitem[Gladman(1993)]{Gladman1993} Gladman, B. 1993, Icarus, 106, 247
\bibitem[Laskar \& Petit(2017)]{LaskarPetit2017} Laskar, J., \& Petit,
A.~C. 2017, A\&A, 605, A72
\bibitem[Li et al.(2022)]{Li2022} Li, Z., et al. 2022, AJ, 164, 163
\bibitem[Qin et al.(2025)]{Qin2025} Qin, Z., Liu, S.-F., \& Ma, B. 2025,
ApJL, doi:10.3847/2041-8213/adea69
\bibitem[Rein \& Liu(2012)]{ReinLiu2012} Rein, H., \& Liu, S.-F. 2012,
A\&A, 537, A128
\bibitem[Rein \& Tamayo(2015)]{ReinTamayo2015} Rein, H., \& Tamayo, D.
2015, MNRAS, 452, 376
\bibitem[Rivera et al.(2010)]{Rivera2010} Rivera, E.~J., Laughlin, G.,
Butler, R.~P., et al. 2010, ApJ, 719, 890
\bibitem[Tamayo et al.(2021)]{Tamayo2021} Tamayo, D., Gilbertson, C., \&
Foreman-Mackey, D. 2021, MNRAS, 501, 4798
\bibitem[Vogt et al.(2014)]{Vogt2014} Vogt, S.~S., et al. 2014, ApJ,
787, 97
\bibitem[Volpi et al.(2019)]{Volpi2019} Volpi, M., Roisin, A., \& Libert,
A.-S. 2019, A\&A, 626, A74
\end{thebibliography}

\end{document}
